ロボットの模倣学習において，視覚的な遮蔽や柔軟物の変形を伴うタスクの遂行には，視覚情報だけでなく触覚情報の活用が不可欠であるが，その触覚情報として従来の高精度な力覚センサ等を用いると，コストがかかり，システムも複雑になる\cite{example}．
本報告では，単純な構造を持つイオン液体センサを触覚とするソフトフィンガーを用い，その低次元な変形情報を模倣学習に統合する手法を提案する．また，7自由度ロボットアームを用いた繊細な接触動作の教示データ収集を容易にするため，重力補償を導入したテレオペレーションシステムを構築し，学習と検証を行った．提案手法の評価として，視覚的遮蔽を伴うシート状柔軟物把持タスクを対象に実験を行った．その結果，触覚情報を統合したモデルは高い成功率を示し，単純なセンサであっても環境変動に対してロバストな動作生成ができることが示された．
